% ! TEX root = ../am2-20-21.tex
\chapter{Ecuații cu derivate parțiale de ordinul doi (EDP 2)}

\section{Clasificare și forma canonică}
\index{EDP!ordinul doi}
O ecuație \emph{cvasiliniară} cu derivate parțiale de ordinul al doilea,
cu două variabile independente, are forma generală:
\[
    A(x, y) z_{xx} + 2B(x, y) z_{xy} + C(x, y) z_{yy} + D(x, y, z, z_x, z_y) = 0,
\]
unde $ A, B, C : \RR^2 \to \RR$ sînt funcții reale, continue pe un deschis
din $ \RR^2 $ (care este domeniul de definiție al ecuației), iar $ D $ este
de asemenea funcție continuă.
\index{EDP!curbe caracteristice}
\begin{definition}\label{def:curbe-car}
    Se numesc \emph{curbe caracteristice} ale ecuației de mai sus curbele care se
    află pe suprafețele integrale ale ecuației, i.e.\ care satisfac
    \emph{ecuația caracteristică}:
    \[
        A(x, y) dy^2 - 2B(x, y) dxdy + C(x, y) dx^2 = 0.
    \]
\end{definition}

Clasificarea ecuațiilor se face în funcție de curbele caracteristice.
Astfel, avem:
\begin{itemize}
    \item $ AC - B^2 < 0 \Rightarrow $ ecuația este de \emph{tip hiperbolic};
    \item $ AC - B^2 = 0 \Rightarrow $ ecuația este de \emph{tip parabolic};
    \item $ AC - B^2 > 0 \Rightarrow $ ecuația este de \emph{tip eliptic}.
\end{itemize}

Exemple foarte des întîlnite provin din fizica matematică:
\begin{itemize}
    \item \textbf{Ecuația coardei vibrante}, care are aceeași formă generală
        cu \textbf{ecuația undelor plane}: \index{EDP!ecuația coardei vibrante} \index{EDP!ecuația undelor plane}
        \[
            u_{xx} - \dfrac{1}{a^2} u_{tt} = 0, \quad a^2 = \dfrac{\rho}{T_0},
        \]
        unde $ \rho $ este densitatea liniară a coardei, iar $ T_0 $ este tensiunea
        la care este supusă coarda în poziția de repaus.

        Ecuația este de tip hiperbolic, după cum se poate verifica imediat.
    \item \textbf{Ecuația căldurii}, cu forma generală: \index{EDP!ecuația căldurii}
        \[
            u_{xx} = \dfrac{1}{a^2} u_t, \quad a^2 = \dfrac{k}{c \rho},
        \]
        unde $ k $ este coeficientul de conductibilitate termică, $ c $ este
        căldura specifică, iar $ \rho $ este densitatea.

        Ecuația are tip parabolic.
    \item \textbf{Ecuația lui Laplace}, cu forma generală: \index{EDP!ecuația Laplace}
        \[
            \Delta u = u_{xx} + u_{yy} = 0,
        \]
        care este o ecuație de tip eliptic.
\end{itemize}

%%%%%%%%%%%%%%%%%%%%%%%%%%%%%%%%%%%%%%%%%%%%%%%%%%%%%%%%%%%%%%%%%%%%%% 

\section{Forma canonică}

Pornind de la ecuația caracteristică, o putem rezolva ca pe o ecuație de gradul
al doilea și obținem, în general, două soluții:
\[
    \dfrac{dy}{dx} = \mu_1(x, y), \quad \dfrac{dy}{dx} = \mu_2(x, y).
\]

Aceste ecuații se numesc \emph{curbele caracteristice} ale ecuației de pornire.

Prin integrarea celor două ecuații, se obțin două familii de curbe în planul
$ XOY $, de forma $ \varphi_1(x, y) = c_1 $ și $ \varphi_2(x, y) = c_2 $, unde
$ c_1 $ și $ c_2 $ sînt constante arbitrare.

Aducerea ecuației de pornire la forma canonică se face pe următoarele cazuri:
\begin{itemize}
    \item Dacă ecuația este \emph{de tip hiperbolic}, se face schimbarea de variabile:
        \[
            \tau = \varphi_1(x, y), \quad \eta = \varphi_2(x, y),
        \]
        iar prima formă canonică a ecuației se obține a fi:
        \[
            z_{\tau\eta} + \Psi_1(\tau, \eta, z, z_\tau, z_\eta) = 0,
        \]
        Putem face și transformarea $ \tau = x + y $ și $ \eta = x - y $, iar a doua
        formă canonică se obține a fi:
        \[
            z_{xx} - z_{yy} + \Phi_1(\tau, \eta, z, z_x, z_y) = 0.
        \]
    \item Dacă ecuația este \emph{de tip parabolic}, o să obținem
        $ \varphi_1 = \varphi_1 = \varphi(x, y) $ și vom face schimbarea de variabilă:
        \[
            \tau = \varphi(x, y), \quad \eta = x.
        \]
        Forma canonică este:
        \[
            z_{\eta\eta} + \Psi_2(\tau, \eta, z, z_\tau, z_\eta) = 0.
        \]
    \item Pentru ecuațiile \emph{de tip eliptic}, funcțiile $ \varphi_1 $ și $ \varphi_2 $
        sînt complex conjugate și putem nota $ \alpha(x, y) = \dr{Re} \varphi_1(x, y) $,
        iar $ \beta(x, y) = \dr{Im}\varphi_1(x, y) $. Schimbarea de variabile este:
        \[
            \tau = \alpha(x, y), \quad \eta = \beta(x, y),
        \]
        iar forma canonică este:
        \[
            z_{\tau\tau} + z_{\eta\eta} + \Psi_3(\tau, \eta, z, z_\tau, z_\eta) = 0.
        \]
\end{itemize}

%%%%%%%%%%%%%%%%%%%%%%%%%%%%%%%%%%%%%%%%%%%%%%%%%%%%%%%%%%%%%%%%%%%%%% 
\section{Cazul coeficienților constanți și $D = 0$}
\label{sec:coef-ct-d0}

Ne ocupăm deocamdată de cazul coeficienților constanți și $ D = 0 $, ecuația prezentîndu-se
în forma generală:
\[
    A z_{xx} + 2B z_{xy} + C z_{yy} = 0, \quad A, B, C \in \RR.
\]

Atunci ecuația diferențială a curbelor caracteristice este:
\[
    Ady^2 - 2Bdxdy + Cdx^2 = 0.
\]
Obținem soluțiile de forma generală:
\[
    \begin{cases}
        dy - \mu_1 dx &= 0 \\
        dy - \mu_2 dx &= 0
    \end{cases} \Rightarrow
    \begin{cases}
        y - \mu_1 x &= c_1 \\
        y - \mu_2 x &= c_2
    \end{cases}, c_1, c_2 \in \RR.
\]

Aducerea la forma canonică se simplifică:
\begin{itemize}
    \item Pentru cazul hiperbolic, substituția este:
        \[
            \begin{cases}
                \tau &= y - \mu_1 x \\
                \eta &= y - \mu_2 x
            \end{cases},
        \]
        iar ecuația devine:
        \[
            z_{\tau\eta} = 0,
        \]
        care are soluția generală $ z = f(\tau) + g(\eta) $, unde $ f $ și $ g $ sînt
        funcții arbitrare. Înlocuim în vechile variabile și obținem:
        \[
            z(x, y) = f(y - \mu_1 x) + g(y - \mu_2 x).
        \]
    \item Pentru cazul parabolic, avem $ \mu_1 = \mu_2 = \dfrac{B}{A} $, iar ecuația
        curbelor devine $ Ady - Bdx = 0 $, care are soluția $ Ay - Bx = c \in \RR $.

        Schimbarea de variabile este $ \tau = Ax - By $, iar $ \eta = x $, care conduce
        la forma canonică:
        \[
            z_{\eta\eta} = 0.
        \]
        Soluția generală este $ z = \eta f(\tau) + g(\tau) $, cu $ f, g $ funcții arbitrare.
        Putem reveni la variabilele anterioare și găsim:
        \[
            z = x f(Ax - By) + g(x).
        \]
    \item Pentru cazul eliptic, forma canonică este chiar ecuația Laplace:
        \[
            z_{\tau\tau} + z_{\eta\eta} = 0,
        \]
        care nu se poate rezolva ușor pe cazul general.
\end{itemize}

%%%%%%%%%%%%%%%%%%%%%%%%%%%%%%%%%%%%%%%%%%%%%%%%%%%%%%%%%%%%%%%%%%%%%%
\section{Exerciții}

1. Aduceți la forma canonică următoarele ecuații:
\begin{enumerate}[(a)]
    \item $ \dfrac{\partial^2 u}{\partial x^2} + 2 \dfrac{\partial^2 u}{\partial x \partial y} - %
        3 \dfrac{\partial^2 u}{\partial y^2} + 2 \dfrac{\partial u}{\partial x} + %
        6 \dfrac{\partial u}{\partial y} = 0 $;
    \item $ 3 \dfrac{\partial^2 u}{\partial x^2} + 7 \dfrac{\partial^2 u}{\partial x\partial y} + %
        2 \dfrac{\partial^2 u}{\partial y^2} = 0 $;
    \item $ 4 \dfrac{\partial^2 u}{\partial x^2} + 4 \dfrac{\partial^2 u}{\partial x \partial y} + %
        \dfrac{\partial^2 u}{\partial y^2} - 2 \dfrac{\partial u}{\partial y} = 0 $;
    \item $ \dfrac{\partial^2 u}{\partial x^2} - 6 \dfrac{\partial^2 u}{\partial x\partial y} + %
        10 \dfrac{\partial^2 u}{\partial y^2} + \dfrac{\partial u}{\partial x} - %
        3 \dfrac{\partial u}{\partial y} = 0 $;
    \item $ 2 \dfrac{\partial^2 u}{\partial x^2} - 7 \dfrac{\partial^2 u}{\partial x \partial y} + %
        3 \dfrac{\partial^2 u}{\partial y^2} = 0 $;
    \item $ y \dfrac{\partial^2 u}{\partial x^2} + (x + y) \dfrac{\partial^2 u}{\partial x \partial y} + %
        x \dfrac{\partial^2 u}{\partial y^2} = 0 $;
    \item $ (1 + x^2) \dfrac{\partial^2 u}{\partial x^2} + %
        (1 + y^2) \dfrac{\partial^2 u}{\partial y^2} + %
        x \dfrac{\partial u}{\partial x} + y \dfrac{\partial u}{\partial y} = 0 $;
    \item $ x^2 \dfrac{\partial^2 u}{\partial x^2} - %
        2xy \cdot \dfrac{\partial^2 u}{\partial x \partial y} + %
        y^2 \cdot \dfrac{\partial^2 u}{\partial y^2} + %
        x \cdot \dfrac{\partial u}{\partial x} + y \cdot \dfrac{\partial u}{\partial y} = 0 $;
\end{enumerate}

\emph{Soluție: (a)} Deoarece avem $ A = 1, B = 1, C = -3 $, rezultă $ B^2 - AC = 4 > 0 $,
deci ecuația este de tip hiperbolic.

Scriem ecuația caracteristică:
\[
    \Big( \dfrac{dy}{dx} \Big)^2 - 2 \dfrac{dy}{dx} - 3 = 0,
\]
care poate fi rezolvată ca o ecuație de gradul al doilea, de unde:
\[
    \begin{cases}
        \dfrac{dy}{dx} = 3 &\Rightarrow y - 3x = c_1 \\
        \dfrac{dy}{dx} = -1 &\Rightarrow y + x = c_2
    \end{cases}.
\]

Cu schimbarea de variabile:
\[
    \begin{cases}
        \tau &= y - 3x \\
        \eta &= y + x
    \end{cases},
\]
funcția căutată $ u(x, y) $ devine $ u(\tau, \eta) $, astfel că toate derivatele parțiale
se calculează acum folosind formula funcțiilor implicite și derivarea funcțiilor compuse.
Toate derivatele parțiale în raport cu $ x $ și $ y $ se calculează, deci, ținînd cont de legătura
cu noile variabile $ \eta $ și $ \tau $. Practic, avem:
\[
    \dfrac{\partial}{\partial x} = \dfrac{\partial}{\partial \eta} \cdot \dfrac{\partial \eta}{\partial x} + %
    \dfrac{\partial}{\partial \tau} \cdot \dfrac{\partial \tau}{\partial x}
\]
și similar pentru $ y $.

Obținem, succesiv:

\newpage
\begin{align*}
    \dfrac{\partial u}{\partial x} &= \dfrac{\partial u}{\partial \tau} \cdot \dfrac{\partial \tau}{\partial x} + %
    \dfrac{\partial u}{\partial \eta} \cdot \dfrac{\partial \eta}{\partial x} \\
                                   &=- 3\dfrac{\partial u}{\partial \tau} + \dfrac{\partial u}{\partial \eta} \\
    \dfrac{\partial u}{\partial y} &= \dfrac{\partial u}{\partial \tau} \cdot \dfrac{\partial \tau}{\partial y} + %
    \dfrac{\partial u}{\partial \eta} \cdot \dfrac{\partial \eta}{\partial y} \\
                                   &=\dfrac{\partial u}{\partial \tau} + \dfrac{\partial u}{\partial \eta} \\
    \dfrac{\partial^2 u}{\partial x^2} &= \dfrac{\partial}{\partial x} \Big( \dfrac{\partial u}{\partial x} \Big)  \\
                                       &= \dfrac{\partial}{\partial x} \Big( \dfrac{\partial u}{\partial \tau} \cdot %
                                       \dfrac{\partial \tau}{\partial x} + \dfrac{\partial u}{\partial \eta} \cdot \dfrac{\partial \eta}{\partial x} \Big) \\
                                       &= \dfrac{\partial}{\partial x} \Big( \dfrac{\partial u}{\partial \tau} \cdot \dfrac{\partial \tau}{\partial x}\Big) + %
                                       \dfrac{\partial}{\partial x} \Big( \dfrac{\partial u}{\partial \eta} \cdot \dfrac{\partial \eta}{\partial x} \Big) \\
                                       &= \dfrac{\partial}{\partial x} \Big( \dfrac{\partial u}{\partial \tau} \Big) \cdot %
                                       \dfrac{\partial \tau}{\partial x} + \dfrac{\partial u}{\partial \tau} \cdot %
                                       \dfrac{\partial^2 \tau}{\partial x^2} + \dfrac{\partial}{\partial x} \Big( \dfrac{\partial u}{\partial \eta} \Big) \cdot %
                                       \dfrac{\partial \eta}{\partial x} + \dfrac{\partial u}{\partial \eta} \cdot
                                       \dfrac{\partial^2 \eta}{\partial x^2} \\
                                       &= \Big[ \dfrac{\partial}{\partial \tau} \Big(\dfrac{\partial u}{\partial \tau}\Big) \cdot %
                                           \dfrac{\partial \tau}{\partial x} + \dfrac{\partial}{\partial \eta} \Big(\dfrac{\partial u}{\partial \tau}\Big) \cdot %
                                       \dfrac{\partial \eta}{\partial x} \Big] \cdot \dfrac{\partial \tau}{\partial x} + \\
                                       &+ \Big[ \dfrac{\partial}{\partial \eta} \Big(\dfrac{\partial u}{\partial \eta}\Big) \cdot %
                                           \dfrac{\partial \eta}{\partial x} + \dfrac{\partial}{\partial \tau} \Big(\dfrac{\partial u}{\partial \eta}\Big) \cdot %
                                       \dfrac{\partial \eta}{\partial x} \Big] \cdot \dfrac{\partial \eta}{\partial x} + \\
                                       &+ \dfrac{\partial^2 \tau}{\partial x^2} \cdot \dfrac{\partial u}{\partial \tau} + %
                                       \dfrac{\partial u}{\partial \eta} \cdot \dfrac{\partial^2 \eta}{\partial x^2} \\
                                       &= \dfrac{\partial^2 u}{\partial \tau^2} \cdot \Big( \dfrac{\partial \tau}{\partial x} \Big)^2 + %
                                       \dfrac{\partial u}{\partial \eta \partial \tau} \cdot \dfrac{\partial \eta}{\partial x} \cdot \dfrac{\partial \tau}{\partial x} + %
                                       \dfrac{\partial^2 u}{\partial^2 \eta} \cdot \Big( \dfrac{\partial \eta}{\partial x} \Big)^2 + \\
                                       &+ \dfrac{\partial^2 u}{\partial \tau\partial \eta} \cdot \dfrac{\partial \tau}{\partial x} \cdot \dfrac{\partial \eta}{\partial x} + %
                                       \dfrac{\partial^2 \tau}{\partial x^2} \cdot \dfrac{\partial u}{\partial \tau} + \dfrac{\partial u}{\partial \eta} \cdot %
                                       \dfrac{\partial^2 \eta}{\partial x^2} \\
                                       &= 9 \dfrac{\partial^2 u}{\partial \tau^2} - %
                                       6 \dfrac{\partial^2 u}{\partial \tau \partial \eta} + %
                                       \dfrac{\partial^2 u}{\partial \eta^2} \\
    \dfrac{\partial^2 u}{\partial x \partial y} &= \dfrac{\partial}{\partial x}\Big( \dfrac{\partial u}{\partial y} \Big)\\
                                                &= -3\dfrac{\partial^2 u}{\partial \tau^2} - %
                                                2 \dfrac{\partial ^2 u}{\partial \tau \partial \eta} + %
                                                \dfrac{\partial^2 u}{\partial \eta^2} \\
    \dfrac{\partial^2 u}{\partial y^2} &= \dfrac{\partial^2 u}{\partial \tau^2} + %
    2 \dfrac{\partial^2 u}{\partial \tau \partial \eta} + %
    \dfrac{\partial^2 u}{\partial \eta^2}.  
\end{align*}

\newpage

Rezultă că, în final, forma canonică este:
\[
    -16 \dfrac{\partial^2 u}{\partial \tau \partial \eta} + 8 \dfrac{\partial u}{\partial \eta} = 0 %
    \Leftrightarrow \dfrac{\partial^2 u}{\partial \tau \partial \eta} - %
    \dfrac{1}{2} \dfrac{\partial u}{\partial \eta} = 0.
\]

\vspace{0.3cm}

2. Determinați soluția ecuației:
\[
    \dfrac{\partial^2 u}{\partial x^2} + 2 \dfrac{\partial^2 u}{\partial x \partial y} - %
    3 \dfrac{\partial^2 u}{\partial y^2} = 0,
\]
care satisface condițiile:
\[
    \begin{cases}
        u(x, 0) &= 3x^2 \\
        \dfrac{\partial u}{\partial y}(x, 0) &= \cos x
    \end{cases}.
\]

\vspace{0.3cm}

3. Rezolvați ecuația:
\[
    3 \dfrac{\partial^2 u}{\partial x^2} + 7 \dfrac{\partial^2 u}{\partial x\partial y} + %
    2 \dfrac{\partial^2 u}{\partial y^2} = 0,
\]
cu condițiile:
\[
    \begin{cases}
        u(x, 0) &= x^3 \\
        \dfrac{\partial u}{\partial y}(x, 0) &= 2x^2
    \end{cases}.
\]

\textit{Indicații pentru 2. și 3.:} Aduceți ecuațiile la forma canonică și
apoi verificați dacă vă aflați în unul dintre cazurile simple din
\S\ref{sec:coef-ct-d0}, care se rezolvă simplu.

